% res_compare.tex -- DRAFT 23 Oct 2026 (thesis result C): our planner against our port of Quan et al.'s method, twin.
% Updated 23 Oct afternoon with the search's real best point (lab batch 17).
% Sources: card docs/experiments/2026-10-23_paper_method_compare.md (critic 22 Oct, Response 22 Oct; deviations section
% 3.1; Results scored 14:01 with tools/c_metrics.py); results/analysis/network_tune_search_2026-10-23.json (the offline
% retune); decisions.md DECIDED 2026-10-23 morning (6) (report as is, no harness fix, the card's seven values accepted)
% and late morning (1) (fly the search's real best point, c_net_portfinal: diss_c 12, k_pay = k_anchor 80, k_ring 40,
% k_slot 18, ki_load 0.5, vel_k_att 8, vel_kp_pos 1, vel_kv 3; J 115, every offline guard passes).
% Registry: R2446 (c_net_smoke_hover_L100, as ported), R2447 (c_net_smoke_linex_v200_L100, as ported), R2448
% (c_pl_linex_v400_L100, our planner), R2468/R2469 (c_net_tune_portbest_*: the search's STAGE-1 point, diss_c 12 and
% vel_kp_pos 1 with the rest at the DEFAULTS, flown first in error as "the best"; corrected 23 Oct ~10:00), R2685-R2690
% (c_net_tune_portfinal_*: the search's best, hover / line 0.2 / line 0.4, two each). Our planner's hover: the B baseline
% b4_hover_L100_ema R2348/R2351, scored by tools/c_metrics.py read-only (card section 4; LAND at 40 s cut the 60 s
% window).
% As-ported numbers are the card's Results (c_metrics, in the scored window): yaw largest change 124.7 / 473.7 deg, tilt
% peak 44.81 / 57.67 deg (the registry's whole-run values, yaw range 230 / 602 deg and tilt 44.8 / 67.7 deg, run past the
% loss). Lost runs: their errors are left out of the cell (card section 5).
% PENDING (card Results): our planner's hover x2 (c_pl_hover_L100), line 0.2 x2 (c_pl_linex_v200_L100) and line 0.4 x1
% more (c_pl_linex_v400_L100), so that every case has both arms at n = 2; they wait for the lab (busy since 14:00). The
% paper's form was never flown in the twin. The claim (3 runs per cell on one code tree) was not flown; Wesley: report
% as is, the bars matter less than that it works.
% Must not say: that Quan et al.'s method cannot hold the ring (the capsize comes from our port's yaw-following formation,
% the locked yaw-only design; the paper's payload is a point with no yaw); "dissipative" for our controller.

\subsection{Our planner against our port of Quan et al.'s method}
\label{sec:res_compare}

This section asks how closely our port of Quan et al.'s method \cite{quan2025swarm} carries the ring compared with
our planner. The method holds a team in formation with a virtual spring--damper network and drives each UAV through
a velocity loop; a member leaving or joining needs no new solver, which is why it is the natural comparison for a
team that changes in flight. Both were flown in the rig twin with four UAVs on plates 0, 3, 6 and 9, 1\,m cables and
the ring at 1.0\,m, in hover and on a 0.9\,m line out and back. Up to the hand-over both arms are our planner (the
creep, the lift and a 5\,s hover). After it, our arm keeps our planner and its trackers with the model integral,
while the port's network drives the trackers' velocity loop with no integral per UAV. The comparison is therefore
between the two methods as flown after the lift, not of the network alone. Every run here is exploratory, one or two
per cell (Table~\ref{tab:compare_twin}), and our planner's hover and its line at 0.2\,m/s are still to fly.

\paragraph{The port.} The port follows the authors' description and simulation code where the rig allows, and
departs from the paper in these ways:
\begin{itemize}
  \item \emph{Nodes:} integrated virtual nodes, sprung to a payload node pinned at the target and to an anchor above
  it, as in the authors' simulation code, where the paper's nodes are a kinematic projection.
  \item \emph{Damping:} against the pinned nodes as if they were still, where the paper's friction is relative to the
  formation's centre.
  \item \emph{The command:} our trackers' velocity loop stands in for the authors' autopilot; the thrust feed-forward
  carries a modelled rod pull, not a measured one; the expected acceleration is not fed forward.
  \item \emph{Rod lengths:} at the retuned gains, the rods the planner measured at the hand-over.
  \item \emph{Yaw:} the formation turns with the ring's measured yaw; the paper's payload is a point with no yaw.
  \item \emph{Takeoff:} our planner creeps and lifts; the network takes over 1.5\,s after the lift tops out, seeded at
  the measured UAVs, with its datum moved onto the measured ring.
  \item \emph{Integrators:} one common height trim for the team, none per UAV.
  \item \emph{Plant:} the twin's rigid rods.
\end{itemize}

\paragraph{As ported, the ring capsizes in yaw.} At its first gains the port lost the ring in the twin, in hover and
on the line at 0.2\,m/s: the ring's yaw swung with a growing amplitude until the ring tipped past the abort limits
(Table~\ref{tab:compare_twin}, Figure~\ref{fig:compare_yaw}). The offline harness had predicted this mode: kicked
0.5$^\circ$ in yaw, the ring's yaw grew about one and a half times a second, then it spun and capsized.

\paragraph{Retuned, the port holds the ring but not its yaw.} The retune was set before the comparison flew and
bounded in time: offline, a grid over the network's damping and the velocity loop's position gain, then a descent over
eight gains from the grid's best point, ranked first on guards (no divergence, ring tilt at most 15$^\circ$, UAV tilt at
most 35$^\circ$, throttle at most 0.78, the ring's yaw within 10$^\circ$ and not growing, the hover no worse) and then on
the line's error at 0.4\,m/s (Table~\ref{tab:compare_offline}). The descent's best point passed every offline guard.
Flown in the twin, twice per case, it held the ring in every run, but not its yaw: the ring's yaw swung through tens of
degrees in hover and on the slow line, and least on the fast line, against the guard's 10$^\circ$, and the ring's tilt
peaked well above our planner's. Offline and the twin disagree on the yaw. The harness hangs the
rods from the UAVs' centres and leaves out the 4\,cm pivot below them, which turns a position mismatch into a torque;
that is the likeliest reason it misses the twin's yaw mode, but it was not tested. The paper's own form of the damping
and feed-forward failed only the tilt guard on the 0.6\,m/s line, a speed this comparison does not fly; it is to fly in
the twin.

\paragraph{On the line, a second late and short of the ends.} On the line at 0.4\,m/s the retuned port ran about a
second behind its schedule and turned short of each end, as the harness had predicted. Its error, the distance across
the line plus the shortfall at the ends, was about twice our planner's (Table~\ref{tab:compare_twin},
Figure~\ref{fig:compare_line}). Our planner also turned short of the ends, by less, and ran a few hundredths of a
second late. The error across the line alone would hide the port's shortfall, which is why the comparison adds the
reach. At 0.2\,m/s the port ran later still and turned closer to the ends; our planner's line at that speed is still to
fly.

\begin{table}[htbp]
\centering
\footnotesize
\setlength{\tabcolsep}{2.6pt}
\caption[Our planner against our port in the twin]{Our planner against our port of Quan et al.'s method in the rig
twin: four UAVs on plates 0, 3, 6 and 9, the 0.86\,kg ring at 1.0\,m on 1\,m cables; $n$ runs each, exploratory (our
planner's hover and its line at 0.2\,m/s, and a second run of its line at 0.4\,m/s, are to fly). Both arms are scored
against the target from just before the hand-over. Error: in hover the ring's 3-D error from the hand-over; on the
line the error across the line and in height (clock-free) plus the reach, the shortfall at the line's ends ($+$ short).
Lag: behind the schedule. Yaw: the ring's largest change of yaw in the window. Search's best: the retune's best point
offline, with the measured rods. $^a$\,The four-UAV hover of the baseline set, on a shorter window (LAND at 40\,s). $^b$\,The ring's yaw grew until
it tipped past the abort limits: aborted 27\,s (hover) and 31\,s (line) after the hand-over; the error of a lost run is
not scored.}
\label{tab:compare_twin}
\begin{tabularx}{\textwidth}{>{\raggedright\arraybackslash}p{1.9cm}>{\raggedright\arraybackslash}p{1.3cm}c>{\raggedright\arraybackslash}Xcccccc}
\toprule
 & & & & Error & Clock-free / & Lag & Ring tilt & Yaw \\
Arm & Case & $n$ & Outcome & (mm) & reach (mm) & (s) & peak ($^\circ$) & ($^\circ$) \\
\midrule
Our planner & hover$^a$ & 2 & held & 1.3--1.4 & -- & -- & 0.21 & -- \\
 & hover & [2] & [to fly] & [\dots] & -- & -- & [\dots] & [\dots] \\
 & line, 0.2\,m/s & [2] & [to fly] & [\dots] & [\dots] & [\dots] & [\dots] & [\dots] \\
 & line, 0.4\,m/s & 1 & held & 46.2 & 1.7 / 44.5 & $+$0.03 & 3.5 & 0.2 \\
\addlinespace
Our port, as ported & hover & 1 & capsized$^b$ & -- & -- & -- & 44.8 & 125 \\
 & line, 0.2\,m/s & 1 & capsized$^b$ & -- & -- & -- & 57.7 & 474 \\
\addlinespace
Our port, the search's best & hover & 2 & held & 16.1--18.7 & -- & -- & 8.4--12.9 & 22--51 \\
 & line, 0.2\,m/s & 2 & held & 41.5--42.1 & 9.1--9.9 / 32.2--32.4 & $+$1.25 & 5.0--8.2 & 29--48 \\
 & line, 0.4\,m/s & 2 & held & 99.3--104.5 & 9.3--9.5 / 89.8--95.1 & $+$1.12--1.13 & 11.2--16.4 & 14 \\
\bottomrule
\end{tabularx}
\end{table}
% R2348/R2351 | c_pl_hover_L100 x2 (to fly) | c_pl_linex_v200_L100 x2 (to fly) | R2448 (+ one more to fly) | R2446
% (hand-over 27.47 s, FLEET_ABORT 54.1 s) | R2447 (abort 58.4 s) | R2685/R2686 (first 5 s 22.4/22.5 mm; height -2.9/-3.1
% mm; quad tilt peak 12.9/11.6 deg; throttle peak 0.635/0.615) | R2687/R2688 (reach mean +31.5/+31.2 mm; time-indexed
% error 170.2/170.0 mm; amplitude ratio 0.930/0.931) | R2689/R2690 (reach mean +94.1/+86.3 mm; time-indexed 279.2/278.3
% mm, ours 33.2; amplitude ratio 0.797/0.800; quad tilt 12.0/12.7 deg; throttle 0.624/0.632; path start after the top
% 4.98 s, so net_hold_carry passes) | R2468 (first 5 s 19.0 mm; yaw first over 10 deg at hand-over + 38.9 s; growth
% x1.14/s against x1.43/s as ported) | R2469 (yaw first over 10 deg at + 33.8 s; time-indexed 261.0 mm; amplitude ratio
% 0.712 vs ours 0.909). R2446/R2447 flew with net_twin_geometry off; R2468/R2469 and R2685-R2690 with it on (its own bar
% run was never flown). After LAND the network's quads drop from the cone (quad tilt 174-176 deg): a note only, the
% landings are not compared.

\begin{table}[htbp]
\centering
\footnotesize
\setlength{\tabcolsep}{3pt}
\caption[The port's retune, offline]{The port's retune in the offline harness, not the twin: the
software-in-the-loop plant at the comparison's geometry (the rig's thrust map, the twin UAVs' offsets, the
motion-capture velocity filter), the network at 10\,Hz and the velocity loop at 50\,Hz, started on the taut
60$^\circ$ cone with the ring 0.5$^\circ$ off in yaw, the rods hung from the UAVs' centres. Each point flies hover, the
line at 0.2, 0.4 and 0.6\,m/s and a margin flight (the line at 0.4\,m/s with the velocity loop's gains 1.25 times
higher); the guards are listed in the text. Paper's form: the damping relative to the target's velocity and the
expected acceleration fed forward. Line: as Table~\ref{tab:compare_twin}. Its numbers are not comparable with the
twin's.}
\label{tab:compare_offline}
\begin{tabularx}{\textwidth}{>{\raggedright\arraybackslash}Xcccl}
\toprule
 & Hover & Line 0.4\,m/s: error & Lag & \\
Form and gains & (mm) & (clock-free / reach) (mm) & (s) & Guards failed \\
\midrule
Port, first gains & \multicolumn{3}{c}{capsized 7--15\,s after the start in every case} & all \\
Paper's form, first gains & \multicolumn{3}{c}{capsized 8--16\,s after the start in every case} & all \\
Port, the grid's point: damping doubled, position gain halved & 8.7 & 151.6 (7.2 / 144.4) & $+$1.01 & margin flight; tilt at 0.6\,m/s \\
Port, the search's best (flown in the twin) & 10.2 & 115.1 (6.9 / 108.3) & $+$1.09 & none \\
Paper's form, the search's best & 8.3 & 12.1 (5.3 / 6.7) & $-$0.01 & tilt at 0.6\,m/s \\
\bottomrule
\end{tabularx}
\end{table}
% network_tune_search_2026-10-23.json: port today's diverged hover 15.4, linex 0.2 12.1, 0.4 12.05, 0.6 14.35, margin
% 7.35 s; paper today's 15.4, 16.1, 11.9, 11.4, 7.6 s (the margin flight is the earliest in both forms). Stage-1 point (12, kp 1): J 151.57, guards
% 'linex_v600 ring tilt 15.3 deg', 'margin_linex_v400 diverged at 17.4 s'. Port best J 115.11 (lag 1.087 s): diss_c 12,
% k_pay = k_anchor 80, k_ring 40, k_slot 18, ki_load 0.5, k_att 8, kp_pos 1, kv 3. Paper best J 12.06 (diss_c 3, k_pay 80,
% k_ring 40, k_slot 18, ki_load 0.5, k_att 16, kp_pos 1.5, kv 3), guard 'linex_v600 ring tilt 20.1 deg'. 54 port and 56
% paper candidates; 12 port candidates pass every guard.

The yaw failure belongs to our port. Its formation turns with the ring's measured yaw, a choice made for the ring of
this work, while the paper's payload is a point with no yaw, so these runs say nothing about Quan et al.'s method as
published. Two further twin tests are to fly: the paper's own form of the damping and feed-forward at its best gains, and
the paper's own geometry, with every cable hooked at the ring's centre.
% Wesley 23 Oct: the grid's-point rows (R2468/R2469, flown by an emit error) and the inferred trade-off sentence cut;
% the paper's form (c_net_tune_paperfinal_*) and the centre-hooked geometry to fly (docs/decisions.md).

\pendingfig{fig:compare_yaw}{The ring's yaw after the hand-over in the rig twin, hover, four UAVs on 1\,m cables: our
port as ported (it capsizes, and the run aborts), our port at the search's best gains (the yaw swings but the ring is
held), and our planner (grey; it does not hand over, so its window starts where the port's would).}{the twin hovers of
both arms, our planner's hover cell included}
% Data: R2446, R2685 (or R2686), our planner's c_pl_hover_L100 (to fly; R2348 meanwhile); run.csv ring quaternion ->
% yaw; window from NET_HANDOVER, ours from the top + 1.5 s. Tool: tools/yaw_excursion.py and tools/c_metrics.py windows ->
% a new thesis_sim_figures function through the thesis-figure skill.

\pendingfig{fig:compare_line}{The line at 0.4\,m/s in the rig twin, four UAVs on 1\,m cables, from the path's start:
(a) the ring's position along the line against the schedule (dotted), our planner (gold) and our port at the search's
best gains (blue; about a second late and short of each end); (b) the ring's yaw over the same window.}{the twin line
runs at 0.4\,m/s of both arms}
% Data: R2448 (ours), R2689 or R2690 (the port's best): run.csv payload_x and the reference; PATH_START from events.csv.
% Tool: tools/c_metrics.py windows -> thesis_sim_figures through the thesis-figure skill. Palette: the port is context,
% so grey may suit it better than a colour; the thesis-figure skill decides.

\paragraph{Answer.} In the rig twin our port of Quan et al.'s method as ported capsizes the ring in yaw. At the best
gains of a bounded retune it holds the ring in every run, but not its yaw, which swings most in hover; on the line at
0.4\,m/s its error is about twice our planner's, from running about a second late and turning short of the ends. Our
planner held the ring level and on its path on the fast line; its hover and its slow line are still to fly. The
failure to hold the yaw is our port's, not the published method's.
